\section{Base-change for local Ext}

\begin{definition}[The module $H(\II,\FF)$]
\label{akcp-1-1-H}
\dcref{1.1}
\source{MR0555258-compactifying-picard:page-0005}
Let $f:X\to S$ be finitely presented and proper, let $\II$ and $\FF$ be
locally finitely presented $\OO_X$-modules, with $\FF$ flat over $S$.  There
are a locally finitely presented $\OO_S$-module $H(\II,\FF)$ and a universal
map $h(\II,\FF):\II\to\FF\otimes f^*H(\II,\FF)$ representing
$M\mapsto\Hom_X(\II,\FF\otimes f^*M)$ on quasi-coherent $\OO_S$-modules.
For every $T\to S$ and quasi-coherent $M$ the Yoneda map
\[
\Hom_T(H(\II,\FF)_T,M)\longrightarrow
\Hom_{X_T}(\II_T,\FF_T\otimes M)
\]
is an isomorphism, and formation of the pair commutes with base change.
\end{definition}

\begin{proof}
The assertion is local on $S$.  After taking $S$ affine, noetherian
approximation descends $(X,\II,\FF)$ to a noetherian base, where the standard
coherent representability theorem for proper morphisms gives $H$ and $h$ and
its arbitrary-base-change compatibility.  Descent back to $S$ proves the
claim.  Tensoring both arguments by an invertible $\OO_S$-module gives the
canonical identification
\[
H(\II\otimes f^*\LL,\FF\otimes f^*\LL)\simeq H(\II,\FF).
\]
The construction is covariant in $\II$, contravariant in $\FF$, and right
exact in either variable.  Applying right exactness to the variable module
gives $H(\II,\FF)\otimes M\simeq H(\II\otimes f^*M,\FF)$.
\end{proof}

\begin{lemma}[Acyclic cover]
\label{akcp-1-2-acyclic-cover}
\dcref{1.2}
\source{MR0555258-compactifying-picard:page-0006}
For every $\OO_X$-module $\II$ there is a surjection $\JJ\twoheadrightarrow\II$
such that, for every $g:Y\to X$ and quasi-coherent $\FF$, $g^*\JJ$ is acyclic
for $\Hom_Y(-,\FF)$.
\end{lemma}
\begin{proof}
Choose affine opens $U$ meeting every stalk of $\II$ and sections generating
those stalks.  The direct sum of the extensions by zero of the corresponding
free modules surjects onto $\II$.  Pullback commutes with this direct sum and
extension by zero.  Since $g^{-1}U$ is affine,
\[
\Ext_Y^q(g^*\JJ,\FF)=\prod_U H^q(g^{-1}U,\FF|_{g^{-1}U})=0\quad(q>0).
\]
\end{proof}

\begin{theorem}[Local freeness criterion]
\label{akcp-1-3-local-free}
\dcref{1.3}
\source{MR0555258-compactifying-picard:page-0006}
\uses{akcp-1-1-H,akcp-1-2-acyclic-cover}
Let $f:X\to S$ be finitely presented and proper and let $\II,\FF$ be locally
finitely presented and $S$-flat.  If
$\Ext^1_{X(s)}(\II(s),\FF(s))=0$ at $s\in S$, then some open retrocompact
neighborhood $U$ of $s$ has $H(\II,\FF)|_U$ locally free of finite rank.
\end{theorem}
\begin{proof}
The assertion is local on $S$ and descends by noetherian approximation, so
take $S=\Spec A$ with $A$ noetherian local and closed point $s$.  Use the
acyclic-cover lemma in an exact sequence
$0\to\KK\to\JJ\to\II\to0$.  Restriction to $X(s)$ remains exact by flatness.
The diagram of Hom and Ext sequences, together with the assumed vanishing and
the adjunction isomorphisms, gives
$\Ext^1_X(\II,\FF\otimes k(s))=0$.  The half-exact functor
$T(M)=\Ext^1_X(\II,\FF\otimes M)$ therefore vanishes on every finite
$A$-module.  Hence $M\mapsto\Hom_A(H(\II,\FF),M)$ is exact, so the finitely
presented module $H(\II,\FF)$ is free at $s$.  Shrinking to the open locus of
local freeness, which is retrocompact after noetherian approximation, proves
the result.
\end{proof}

\begin{lemma}[Free presentation on an affine morphism]
\label{akcp-1-4-free-presentation}
\dcref{1.4}
\source{MR0555258-compactifying-picard:page-0007}
If $f:X\to S$ is finitely presented and affine and $\II$ is finitely presented
and $S$-flat, there is an exact sequence
$0\to\KK\to\JJ\to\II\to0$ with $\KK,\JJ$ finitely presented and $\JJ$
free.
\end{lemma}
\begin{proof}
Descend to a noetherian affine approximation.  There a finite presentation of
the module has a finite free middle term; flatness preserves exactness after
pullback to $S$.
\end{proof}

\begin{lemma}[Ext and inverse limits]
\label{akcp-1-5-ext-limits}
\dcref{1.5}
\source{MR0555258-compactifying-picard:page-0007}
Let $X_0\to S_0$ be finitely presented affine, $S=\varprojlim S_\lambda$ a
projective limit, and let $\II_0,\FF_0$ be finitely presented modules with
$\II_0$ $S_0$-flat for $q>1$ and $X_0$ $S_0$-flat for $q>2$.  For the induced
objects on $X$, the canonical map
\[
\varinjlim_\lambda\Ext^q_{X_\lambda}(\II_\lambda,\FF_\lambda)
\;\xrightarrow{\sim}\;\Ext^q_X(\II,\FF)
\]
is an isomorphism.
\end{lemma}
\begin{proof}
For $q=0$ this is finite-presentation Hom and limit compatibility.  For $q=1$
use the free presentation of \cref{akcp-1-4-free-presentation}; the resulting
Hom diagrams have exact rows because the free term is affine-acyclic, and the
$q=0$ case identifies the limits.  For $q>1$ use the same presentation and
induct on $q$; flatness of $X_0$ makes the free term flat, hence the kernel is
flat and the induction applies.  Homotopy of two free presentations shows that
the dotted maps are canonical.
\end{proof}

\begin{lemma}[Affine change of rings]
\label{akcp-1-6-change-rings}
\dcref{1.6}
\source{MR0555258-compactifying-picard:page-0009}
For a finitely presented $A$-algebra $B$ and finitely presented $B$-modules
$M,N$, if $M$ is $A$-flat for $q\ge1$ and $B$ is $A$-flat for $q\ge2$, then
\[
\Ext_B^q(M,N)\simeq\Ext_A^q(\widetilde M,\widetilde N),
\]
canonically, where the tildes denote the underlying $A$-modules.
\end{lemma}
\begin{proof}
The case $q=0$ is the finite-presentation Hom adjunction.  A finite free
presentation of $M$ and induction on $q$ give the higher cases exactly as in
\cref{akcp-1-5-ext-limits}.
\end{proof}

\begin{lemma}[Adjunction for local Ext]
\label{akcp-1-7-ext-adjunction}
\dcref{1.7}
\source{MR0555258-compactifying-picard:page-0009}
\uses{akcp-1-4-free-presentation,akcp-1-5-ext-limits}
For $f:X\to S$ finitely presented, $\II,\FF$ locally finitely presented, and
$g:T\to S$, there is a functorial adjunction isomorphism
\[
\Ext^q_X\bigl(\II,(1\times g)_*(\FF_T\otimes M)\bigr)
\simeq (1\times g)_*\Ext^q_{X_T}(\II_T,\FF_T\otimes M),
\]
provided $\II$ is $S$-flat for $q\ge1$ and $f$ is flat for $q\ge2$.
\end{lemma}
\begin{proof}
For $q=0$ this is ordinary adjunction.  Resolve $\II$ by a finite free
presentation locally and use induction, with the preceding change-of-rings and
limit compatibilities; all constructions commute with further base changes.
\end{proof}

\begin{definition}[Base-change map]
\label{akcp-1-8-base-change-map}
\dcref{1.8}
\source{MR0555258-compactifying-picard:page-0009}
\uses{akcp-1-7-ext-adjunction}
Under the hypotheses of \cref{akcp-1-7-ext-adjunction}, the canonical map
$\FF\to(1\times g)_*(1\times g)^*\FF$ and tensor adjunction define
\[
b^q(M):\Ext_X^q(\II,\FF)\otimes M\longrightarrow
\Ext_{X_T}^q(\II_T,\FF_T\otimes M).
\]
It is compatible with restrictions, localizations, further base change and the
limits above; if $g$ is flat it is an isomorphism.
\end{definition}

\begin{theorem}[Property of exchange]
\label{akcp-1-9-exchange}
\dcref{1.9}
\source{MR0555258-compactifying-picard:page-0010}
\uses{akcp-1-8-base-change-map}
Let $\FF$ be $S$-flat and assume the flatness hypotheses of
\cref{akcp-1-8-base-change-map}.  If the fiber map
$b^q(k(s))$ is surjective at $x\in X(s)$, then for every $g:T\to S$ and
quasi-coherent $M$, $b^q(M)$ is an isomorphism at every point over $x$.
Moreover the following are equivalent: (1) $b^{q-1}(k(s))$ is surjective at
$x$; (2) all $b^{q-1}(M)$ are isomorphisms over $x$ after arbitrary base change;
(3) $\Ext_X^{q-1}(\II,\FF)$ is $S$-flat at $x$.
\end{theorem}
\begin{proof}
Approximate by noetherian schemes and descend surjectivity to one stage using
finite generation.  Localize at $s,x,t$ and put
$R(N)=\Ext^q_{X_t}(\II_t,\FF_t\otimes N)$.  Surjectivity on the residue
field and Nakayama's lemma give $R(\OO_{S,s})\otimes N\twoheadrightarrow R(N)$;
adjunction identifies this map with $b^q(M)$, proving (i).  For (ii), tensor an
arbitrary short exact sequence and compare its long Ext sequences.  If the
base-change maps are isomorphisms, the tensor map is injective and the Ext
module is flat.  Conversely flatness makes the connecting map injective; use
$M=\OO_S$, $M''=k(s)$ to obtain surjectivity on the fiber.\end{proof}

\begin{theorem}[Open vanishing and exchange loci]
\label{akcp-1-10-open-ext}
\dcref{1.10}
\source{MR0555258-compactifying-picard:page-0012}
\uses{akcp-1-9-exchange}
With $f,\II,\FF$ as above, (i) the locus where
$\Ext^q_{X(s)}(\II(s),\FF(s))=0$ is open and retrocompact, and over it all
base-changed $\Ext^q(\II,\FF\otimes M)$ vanish.  (ii) For fixed $c$, the locus
where the fiber Ext groups in degrees $c+1,c-1$ vanish is open and
retrocompact; $\Ext^c_X(\II,\FF)$ is finitely presented and flat there and all
$b^c(M)$ are isomorphisms.  (iii) If $N\mapsto\Ext^q_X(\II,\FF\otimes N)$ is
right exact, all $b^q(M)$ are isomorphisms; if it is exact, the same holds in
degree $q-1$ and $\Ext^q_X(\II,\FF)$ is $S$-flat.
\end{theorem}
\begin{proof}
The vanishing locus upstairs is open and retrocompact by the local Ext
semicontinuity theorem; properness carries the complement of its fiberwise
image to a closed set.  Apply \cref{akcp-1-9-exchange} to obtain (i).  Apply it
in degrees $c+1$ and $c$ to get (ii), with finite presentation from the standard
local Ext theorem.  For (iii), pull back to $\Spec k(s)$ and use the exact
sequence $0\to\OO_S\to g_*\OO_{k(s)}\to\Coker\to0$; right exactness gives
surjectivity on each fiber, so exchange applies.  Exactness implies right
exactness in adjacent degree and then flatness by the equivalence in
\cref{akcp-1-9-exchange}.
\end{proof}
